
\section{\trArzelaAscoliTheorem}
\label{sec:appC-arzela-ascoli}

让我们回忆这条非常有用的定理：它刻画了\trCompact{}空间上的连续函数空间的相对紧子集，我们在若干个紧性结论中都要用到它。关于这条定理及有关文献，参见 [17]。

\begin{theorem}[\trArzela--\trAscoli]\label{thm:arzela-ascoli-compactness}\BookIndexDefinition{trArzelaAscoliTheorem}
设 $(E,d)$ 是一个\trCompact{}\trMetric{}空间，又设 $\scr{F}$ 是由连续函数 $f : E \to \RR$ 组成的一个\trBounded{}族，诸 $f$ 在 $E$ 上是\trUniform{}\trEquicontinuity{}的，也就是说，满足
\[
\forall \varepsilon > 0,\ \exists \delta > 0\ \text{使得}\ \forall f \in \scr{F},\quad d(x,y) < \delta \Longrightarrow |f(x) - f(y)| < \varepsilon.
\]
那么 $\scr{F}$ 在 $E$ 上的连续函数空间中是相对紧的。
\end{theorem}

换言之，从 $\scr{F}$ 的元素组成的任何一个序列中，我们都能抽取出一个在 $E$ 上\trUniformConvergence{}的序列（\trConvergence{}到一个连续函数）。

\begin{remark}\label{thm:appC-equi-lipschitz}
在本书中我们应用这条定理的所有情形里，函数族都是\trEquicontinuity{}的，原因在于它是“等度\trLipschitz”的。
\end{remark}
